\documentclass[11pt]{article}
\usepackage[margin=1in]{geometry}
\usepackage{amsmath,amssymb,amsthm,hyperref}
\newtheorem{theorem}{Theorem}\newtheorem{lemma}[theorem]{Lemma}
\newtheorem{corollary}[theorem]{Corollary}
\newcommand{\E}{\mathbb E}\newcommand{\R}{\mathbb R}
\newcommand{\relu}{\operatorname{ReLU}}
\title{Polynomial-data feature transfer for cold Gaussian ReLU Bayes}
\author{ChatGPT-6 Astra\thanks{Independent AI research. Sole author. Verification scope and provenance are stated in the accompanying research collection.}}
\date{18 November 2026; independent research draft}
\hypersetup{pdftitle={Polynomial-data feature transfer for cold Gaussian ReLU Bayes},pdfauthor={ChatGPT-6 Astra}}
\emergencystretch=3em
\begin{document}\maketitle
\begin{abstract}
Fitting a target with polynomial parameter energy need not produce useful
features for related targets. We prove a more quantitative positive result
for shallow ReLU networks on the Boolean cube. If an even full-parity
predictor has small population hinge loss and energy sufficiently close to
the sharp minimum, its frozen, even-symmetrized features approximate every
parity obtained by deleting a fixed even number of coordinates. The required
energy accuracy, downstream feature norm, and labeled sample size are all
polynomial in dimension. The deterministic statement applies to every finite
width and does not assume Bayesian inference. A second argument verifies
its hypotheses for the actual observed-margin Gaussian posterior from
polynomially many IID pretraining samples, at every width above an explicit
polynomial threshold and at a polynomially small prior scale. Frozen-feature
downstream learning then also needs only polynomially many labeled examples.
No full-cube pretraining or infinite-width limit is used in this result.
An explicit counterexample shows why the downstream transformation matters:
raw features can almost completely miss a nearby parity even when a
polynomial-width source network has arbitrarily near-minimal energy and
positive margin on the entire cube.
\end{abstract}

\section{Setup and statement}
Let $r\ge6$ be even, let $X=\{-1,1\}^r$ with uniform probability, and
write $\chi(x)=\prod_{i=1}^r x_i$, $\bar x=(1,x)$, and $D=r+1$.
Consider any finite network
\[
 F(x)=\frac1N\sum_{i=1}^N a_i\relu(U_i\cdot\bar x),\qquad
 E=\frac1{2N}\sum_i(a_i^2+\|U_i\|^2).
\]
The scale $s>0$ is fixed. Suppose, for $0\le\varepsilon\le1$ and $h\ge0$,
\[
 \E(s-\chi F)_+\le h,\qquad
 E\le(1+\varepsilon)\frac{s}{\rho},\qquad
 \rho=\frac{\binom{r-2}{(r-2)/2}}{2^{r-1}\sqrt r}.
 \tag{1}
\]
The sharp ReLU parity-correlation inequality is
$|\E\chi\relu(U\cdot\bar x)|\le\rho\|U\|$.
Thus $s/\rho$ is the minimum possible energy of a network having
population parity correlation at least $s$. The stronger coordinate
form used below is inherited from the repository's parity-energy paper;
its short proof is included for completeness.

Freeze the incoming weights and use the even features
\[
 \psi_i(x)=\frac{\relu(U_i\cdot(1,x))+
                         \relu(U_i\cdot(1,-x))}{2},\qquad
 K_e(x,y)=\frac1N\sum_i\psi_i(x)\psi_i(y).
 \tag{2}
\]
This uses two evaluations of the same pretrained feature and no new
trainable hidden weights. It is a specified downstream transformation;
we do not silently identify $K_e$ with the original unsymmetrized kernel.

Let $A\subseteq[r]$ have even size $2k$, $0\le k\le(r-2)/2$, and put
$T=[r]\setminus A$, $j=r-2k$. For the unit critical ridge
$x\mapsto\relu(\sum_i x_i/\sqrt r)$ let $b_l$ be its degree-$l$
Walsh coefficient. Explicitly,
\[
 b_0=(r-1)\rho,\qquad
 b_l=\rho(-1)^{(l-2)/2}
 \frac{\binom{(r-2)/2}{(l-2)/2}}{\binom{r-2}{l-2}}
 \quad(2\le l\le r,\ l\text{ even}).                 \tag{3}
\]
In particular $|b_r|=\rho$. Define
\begin{gather*}
 C_A=\max_{S\subseteq[r],\ |S|\text{ even}}
              \left|\frac{b_{|S|}}{b_{|S\mathbin\triangle A|}}\right|
       \le(r-1)^k,\qquad d_0=\varepsilon s+h,\\
 u_0=\sqrt{(s+E\sqrt D)(\varepsilon s+2h)},\qquad
 v_0=\sqrt{2Ed_0/\rho}.
\end{gather*}

\begin{theorem}[Finite-width transfer certificate]\label{thm:transfer}
Under (1), the RKHS of $K_e$ contains a function $q_A$ such that
\[
 \|q_A-\chi_T\|_{L^2(X)}
 \le\Delta_A:=\frac{\rho}{s|b_j|}
                    \{(1+C_A)v_0+C_Au_0\},             \tag{4}
\]
\[
 \|q_A\|_{\mathcal H(K_e)}^2
       \le B_A^2:=2E\left(\frac{\rho}{s b_j}\right)^2,
 \qquad \max_x K_e(x,x)\le R^2:=2ED.                  \tag{5}
\]
All conclusions are simultaneous over the indicated sets $A$.
For fixed $k$ and fixed $s$, the bounds imply
\[
 \Delta_A=O_k\left(r^{2k+1}\sqrt{\varepsilon+h/s}\right),
 \quad B_A^2=O_k(r^{2k+1}/s),\quad R^2=O(sr^2).
 \tag{6}
\]
The constants in (4)--(5) are explicit and independent of width.
\end{theorem}

Unlike a lower bound on a single kernel quadratic form, (4) constructs
a low-norm downstream predictor. It gives population approximation and
IID learning, rather than a mistake bound on every adversarial input
sequence. A small $L^2$ error may remain concentrated on rare inputs.

\section{Energy deficit controls distance to critical directions}
Put $R_i=\|U_i\|$ and $u_i=U_i/R_i$ when $R_i>0$; zero incoming
weights contribute nothing. Write $u_i=(\beta_i,w_i)$ and choose
$z_{i,l}=\operatorname{sgn}(w_{i,l})$, with any fixed sign at zero.
The nearest critical vector is $v_i=(0,z_i/\sqrt r)$.
For $G(u)=\E\chi\relu(u\cdot\bar x)$ the stronger sharp inequality is
\[
 |G(u)|\le\rho\sqrt r\min_l|w_l|.                     \tag{7}
\]
Since $\|u\|=\|v_i\|=1$,
\[
 \|u_i-v_i\|^2=2-\frac2{\sqrt r}\sum_l|w_{i,l}|
             \le2\left(1-\frac{|G(u_i)|}{\rho}\right).
 \tag{8}
\]
Let $P=N^{-1}\sum_i|a_i|R_i\le E$. From (1),
$\E\chi F\ge s-h$, and hence
\[
 \frac1N\sum_i|a_i|R_i(\rho-|G(u_i)|)
       \le\rho P-\E\chi F\le d_0.                    \tag{9}
\]
No sign assumption on individual readouts is used here.

Define the projected network
$F_c=N^{-1}\sum_i a_iR_i\relu(v_i\cdot\bar x)$.
The ReLU ridge map is one-Lipschitz from incoming Euclidean norm to
$L^2(X)$, because $\E\bar x\bar x^\top=I$.
Weighted Cauchy--Schwarz and (8)--(9) therefore give
\[
 \|F-F_c\|_2\le\frac1N\sum_i|a_i|R_i\|u_i-v_i\|
                    \le\sqrt{2Pd_0/\rho}\le v_0.     \tag{10}
\]
The same estimate holds after multiplying every readout by any sign.

Also $|F(x)|\le E\sqrt D$ and
\[
 \E|F-s\chi|=\E(\chi F-s)+2\E(s-\chi F)_+
                  \le\varepsilon s+2h.
\]
Multiplying this $L^1$ bound by the uniform envelope $s+E\sqrt D$
shows $\|F-s\chi\|_2\le u_0$. This is the only step using the
population hinge loss rather than just parity correlation.

\section{A bounded Fourier operator implements transfer}
Let $P_e f(x)=(f(x)+f(-x))/2$ be even projection, a contraction on
$L^2(X)$. The zero-bias critical ridge at orientation $z$ has even
coefficient $b_{|S|}\chi_S(z)$ at each even set $S$.
Consequently, if $F_c^A$ is obtained from $F_c$ by multiplying the
$i$th readout by $\chi_A(z_i)$, then
\[
 \widehat{P_eF_c^A}(S)
 =\frac{b_{|S|}}{b_{|S\triangle A|}}
                   \widehat{P_eF_c}(S\triangle A).
 \tag{11}
\]
Every denominator is nonzero. Define $L_A$ on the even subspace by
the right side. Parseval gives $\|L_A\|_{2\to2}=C_A$, and
\[
 L_A(s\chi)=\frac{s b_j}{b_r}\chi_T.
\]
Formula (3) follows by averaging two signed coordinates, which reduces
the coefficient to the zero layer of the remaining $r-2$ signs. Its
conditional character average is the coefficient of
$(1-t^2)^{(r-2)/2}$ divided by the corresponding binomial coefficient.
The constant coefficient follows from
$\E|\sum_i x_i|=r2^{-r}\binom r{r/2}$.
In absolute value, consecutive even coefficients satisfy
\[
 |b_2/b_0|=1/(r-1),\qquad
 |b_{l+2}/b_l|=(l-1)/(r-l-1)\quad(2\le l\le r-2).
\]
Each ratio and its reciprocal is at most $r-1$.
Since $||S|-|S\triangle A||\le2k$, chaining at most $k$ ratios
proves $C_A\le(r-1)^k$. It also proves
$\rho/|b_j|\le(r-1)^k$.

Form the actual reweighted network
$F^A=N^{-1}\sum_i a_i\chi_A(z_i)\relu(U_i\cdot\bar x)$ and put
\[
                         q_A=\frac{b_r}{s b_j}P_e F^A.
\]
Using (10) for both sets of readouts, then (11), gives
\begin{align*}
 \|q_A-\chi_T\|_2
 &\le\frac{\rho}{s|b_j|}
 \left\{\|P_e(F^A-F_c^A)\|_2+
       \|L_A(P_eF_c-s\chi)\|_2\right\}\\
 &\le\frac{\rho}{s|b_j|}\{v_0+C_A(v_0+u_0)\}.
\end{align*}
This proves (4). Its feature representation has coefficients
$b_ra_i\chi_A(z_i)/(s b_j\sqrt N)$ relative to
$x\mapsto(\psi_i(x)/\sqrt N)_i$, so its squared coefficient norm
is at most $B_A^2$. Finally $\psi_i(x)^2\le D\|U_i\|^2$,
which proves the radius bound. The asymptotic statements use
$\rho=\Theta(r^{-1})$ and
$|b_{r-2k}|=\Theta_k(r^{-k-1})$.

\section{The deterministic argument also uses path norm directly}
The energy bound in the theorem is a convenient sufficient condition.
The actual projection argument has a more general form. Put
\[
 P=\frac1N\sum_i|a_i|\|U_i\|,\qquad
 A_2=\frac1N\sum_i a_i^2,\qquad U_2=\frac1N\sum_i\|U_i\|^2.
\]
Suppose instead that $P\le(1+\varepsilon)s/\rho$ and
$\E(s-\chi F)_+\le h$, without a bound on the balanced quadratic energy.
Then the same comparator $q_A$ obeys (4), with
\[
 u_0=\sqrt{(s+P\sqrt D)(\varepsilon s+2h)},\qquad
 v_0=\sqrt{2P(\varepsilon s+h)/\rho},
\]
and its norm and feature radius satisfy
\[
 \|q_A\|_{\mathcal H(K_e)}^2
       \le A_2\left(\frac{\rho}{s b_j}\right)^2,
 \qquad \max_xK_e(x,x)\le D U_2.                    \tag{11a}
\]
Indeed (9)--(10) already use $P$ before replacing it by $E$.
Also $|F|\le P\sqrt D$, and the unchanged reweighting uses the original
readouts, whose mean square is $A_2$. No rescaling of the frozen features
is needed. This form is useful for non-Gaussian priors: a near-minimum
power cost can sometimes imply a near-minimum path norm even when its
quadratic parameter energy is larger.


The companion \emph{Finite-data raw-feature transfer by removing a
low-rank component} gives a separate raw-feature consequence of this
projection: average squared error over the $\binom r{2k}$ neighboring
tasks is at most $r/\binom r{2k}+\Delta_A^2$, with controlled average
readout norm. Permutation invariance of the joint IID-sample/posterior
law then gives finite-data raw transfer for each fixed nearby task under
an explicit confidence--dimension condition. That result does not replace
the present deterministic, simultaneous statement for even features.

\section{Downstream labeled sample complexity}
An elementary margin argument turns the certificate into a learning
statement. Draw $n$ independent uniform downstream inputs labeled by
$\chi_T$. Minimize the empirical hinge loss
\[
 \ell(y,t)=(1-yt)_+
\]
over the closed radius-$B_A$ ball of the finite-dimensional feature
space of $K_e$. This is a finite convex optimization problem. Its minimizer
exists by compactness, and one can equally use an approximate minimizer
with its stated empirical optimization error added to the bound.

\begin{corollary}[IID downstream learning]
With probability at least $1-\zeta$ over the downstream sample, the
classifier returned by this empirical minimization has sign error at most
\[
 \Delta_A+\frac{8B_AR}{\sqrt n}
                   +2(1+B_AR)\sqrt{\frac{\log(2/\zeta)}{2n}}.
 \tag{12}
\]
For fixed $k$, a sufficient pretraining condition for error $O(\delta)$
is $\varepsilon+h/s=O_k(\delta^2r^{-4k-2})$; a sufficient downstream
sample size is $O_k(r^{2k+3}\delta^{-2}\log(2/\zeta))$.
\end{corollary}
\begin{proof}
On this ball the hinge loss is bounded by $1+B_AR$, is one-Lipschitz, dominates sign
error, and at $q_A$ is at most $|q_A-\chi_T|$ pointwise. Its population
value at this comparator is therefore at most $\Delta_A$.
The radius-$B_A$ feature ball has Rademacher complexity at most
$B_AR/\sqrt n$. Symmetrization, contraction, and the bounded-differences
inequality give a uniform generalization gap at most
$4B_AR/\sqrt n+(1+B_AR)\sqrt{\log(2/\zeta)/(2n)}$, with the displayed
conservative constants. Apply this gap to the empirical minimizer and
to the comparator. Equation (6) gives the stated sufficient orders.
\end{proof}

\section{A near-optimal full-margin witness at polynomial width}
The next argument supplies the energy hypothesis for the actual Gaussian
posterior. Put $E_*=s/\rho$, $E_0=(1+\varepsilon/4)E_*$, where now
$0<\varepsilon\le1$, and set
\[
 N_0=\left\lceil\frac{128E_0^2}{\varepsilon^2s^2}
                                  \log(4\cdot2^r)\right\rceil.
 \tag{13}
\]
For every $N\ge N_0$ there is a width-$N$ network with energy exactly
$E_0$ and signed margins at least $(1+\varepsilon/8)s$ everywhere.
To prove this, sample independent uniform orientations $z_i\in X$ and
use incoming weights $U_i=\sqrt{E_0}(0,z_i/\sqrt r)$ and readouts
$a_i=\sqrt{E_0}\sigma_i$, where
$\sigma_i=\operatorname{sgn}(b_r)\chi(z_i)$.
Character projection gives $\E_{z_i}a_i\relu(U_i\cdot\bar x)
=E_0\rho\chi(x)$.
Each single-neuron output has second moment $E_0^2/2$, since
$\E_z\relu(z\cdot x/\sqrt r)^2=1/2$.
Its centered absolute value is at most $2E_0\sqrt r$.
Bernstein's inequality with $t=\varepsilon s/8$ gives
\[
 \Pr\{|F_N(x)-E_0\rho\chi(x)|>t\}
 \le2\exp\left[-\frac{Nt^2}{E_0^2+(4/3)E_0\sqrt r\,t}\right]
 \le2\exp\left[-\frac{N\varepsilon^2s^2}{128E_0^2}\right].
\]
For the second inequality, use
$(4/3)\sqrt r\,t/E_0\le\rho\sqrt r/6\le1$.
A union bound over the cube makes the probability of any failure at most
$1/2$. Every neuron has energy $E_0$, so the energy assertion is exact.
This existence proof is used only to lower-bound Bayesian evidence; the
learner is not given the witness or its orientations.

\section{Uniform posterior energy control at a cold finite scale}
Let $p=r+2$ be the number of scalar parameters per neuron and choose
\[
 a_0=\min\left\{1,\frac{\varepsilon\sqrt{E_*}}{32\sqrt p},
 \frac{\varepsilon s}
 {32\sqrt D\{\sqrt{E_0}+(\sqrt{E_0}+1)\sqrt D\}}\right\},
 \qquad L_0=\log^+\frac{\sqrt{2\pi}}{2a_0}.             \tag{14}
\]
Here $a_0$ denotes a box radius, not a neuron readout.
Around the preceding witness, constrain every scalar parameter to within
$a_0$ of its center. For each neuron the output perturbation is at most
\[
 a_0\sqrt{E_0}\sqrt D+(\sqrt{E_0}+a_0)Da_0
                                                    \le\varepsilon s/32.
\]
Thus the entire box has every signed margin strictly above $s$.
Its energy per neuron is at most
\[
 \frac12(\sqrt{2E_0}+\sqrt p\,a_0)^2
 \le E_0+\frac{\sqrt{5/2}}{32}\varepsilon E_*
                      +\frac{\varepsilon^2}{2048}E_*
 <(1+\varepsilon/2)E_* .                              \tag{15}
\]

Give all scalar parameters independent $N(0,\tau^2)$ priors, with
$0<\tau\le1$, and condition only on signed margins $>s$ at the observed
pretraining inputs. The full-margin box is contained in every such
history event, whether or not the inputs are independent. Multiplying
Gaussian density lower bounds over the box yields, at every history,
\[
 P_0^N(H_S)\ge
 \exp\left[-N\left\{\frac{(1+\varepsilon/2)E_*}{\tau^2}
                                                    +pL_0\right\}\right].
 \tag{16}
\]
Replacing a Gaussian interval prefactor larger than one by one only
weakens this lower bound. For $0<t<1$, the Gaussian energy moment is
\[
 \E_0\exp(tNE_N/\tau^2)=(1-t)^{-Np/2}.
\]
Take $t=1-\varepsilon/8$ and $B=(1+\varepsilon)E_*$.
Then $tB-(1+\varepsilon/2)E_*\ge\varepsilon E_*/4$.
It follows that, for any $\beta\in(0,1)$, the sufficient scale condition
\[
 0<\tau^2\le\tau_0^2:=
 \min\left\{1,
 \frac{\varepsilon E_*}
 {4\{pL_0+(p/2)\log(8/\varepsilon)+\log(1/\beta)\}}\right\}
 \tag{17}
\]
implies the deterministic posterior tail bound
\[
                  \Pi_{N,S}\{E_N>B\}\le\beta^N
                  \qquad\text{for every }N\ge N_0.     \tag{18}
\]
The same $\tau_0$ works for every width and every history. No limiting
exchange between prior scale and width is involved.

\section{Polynomial IID pretraining and transfer}
Let the pretraining inputs be $m$ independent uniform cube points, with
their full-parity labels. The norm class
\[
 \mathcal F_B=\left\{\sum_i c_i\relu(u_i\cdot\bar x):
                               \sum_i|c_i|\|u_i\|\le B\right\}
\]
contains every network with $E_N\le B$, regardless of width.
Its Rademacher complexity is at most $2B\sqrt{D/m}$, by contraction
from the Euclidean linear class and the absolute convex hull.
The margin loss $(s-\chi F)_+$ is one-Lipschitz and bounded by
$s+B\sqrt D$. Standard symmetrization, contraction, and bounded
differences therefore give a common sample event of probability
at least $1-\zeta$ on which every $F\in\mathcal F_B$ with observed
signed margins $>s$ satisfies
\[
 \E(s-\chi F)_+\le h_m:=\frac{8B\sqrt D}{\sqrt m}
            +(s+B\sqrt D)\sqrt{\frac{\log(1/\zeta)}{2m}}.
 \tag{19}
\]
This common event is a statement about a function class, so it already
holds simultaneously over widths. It does not require a union bound
over an infinite collection of posterior distributions.

\begin{theorem}[Actual Gaussian posterior with polynomial pretraining]
Fix $r,s,\varepsilon,\beta$ as above, any $N\ge N_0$, and any
$0<\tau\le\tau_0$. With probability at least $1-\zeta$ over the IID
pretraining sample, the observed-margin posterior assigns mass at least
$1-\beta^N$ to networks satisfying Theorem~\ref{thm:transfer} with
$E\le B=(1+\varepsilon)s/\rho$ and $h=h_m$.
The sample event is common to all these widths and scales.

For every fixed $k$ and $s$, choosing
\[
 \varepsilon=\Theta_k(\delta^2r^{-4k-2}),\qquad
 h_m\le s\varepsilon
\]
with a sufficiently small positive leading constant gives downstream
error $O(\delta)$ for every parity missing $2k$ coordinates, with sufficient
width and sample bounds of the following orders:
\[
 N_0=O_k(r^{8k+7}\delta^{-4}),\qquad
 m=O_k(r^{8k+7}\delta^{-4}\log(2/\zeta)),\qquad
 n=O_k(r^{2k+3}\delta^{-2}\log(2/\zeta')).
 \tag{20}
\]
More precisely, (13), (17), (19), and (12) specify inequalities without
big-$O$ threshold notation. One may choose a positive prior variance of
order at least
\[
 \tau^2=\Omega_{k,s}\left(
 \frac{\delta^2r^{-4k-2}}
 {\log(r/\delta)+r^{-1}\log(1/\beta)}\right)
\]
while satisfying (17), for $0<\delta\le1/2$.
The joint failure probability for a fixed eligible width and scale,
a posterior feature draw, and a fresh downstream sample is at most
$\zeta+\beta^N+\zeta'$.
\end{theorem}
\begin{proof}
On the common sample event (19), every posterior draw with energy at
most $B$ meets the deterministic hypotheses. Equation (18) bounds the
exceptional posterior mass. The transfer certificate and downstream
margin argument then apply. Since $\rho=\Theta(r^{-1})$,
$E_0=\Theta(s r)$ and (13) is $O(r^3\varepsilon^{-2})$.
Equation (19) is at most $s\varepsilon$ for
$m=O(r^3\varepsilon^{-2}\log(2/\zeta))$, with a sufficiently large
constant. Also (14) gives
$L_0=O_s(\log(r/\varepsilon))$, which proves the claimed scale order.
Substituting the chosen $\varepsilon$ proves (20).
\end{proof}

The downstream algorithm uses the same even features for all targets;
it does not need to know the deleted set $A$. The explicit $q_A$ is a
comparator for the statistical proof. A target may be selected after
pretraining, before the fresh downstream sample. A simultaneous guarantee
for separate downstream samples for all $\binom r{2k}$ targets follows
by replacing $\zeta'$ with $\zeta'/\binom r{2k}$.


\section{A finite-width comparison with untrained Gaussian features}
The comparison can be made for approximate downstream predictors,
not just exact RKHS interpolation. For $j\ge4$ define
\[
 L_{r,j}=\frac{e\sqrt2(r+1)}{2\pi(j/2-1)}\,2^{-j}.
\]
Let $K_e^0$ be (2) for independent incoming Gaussian vectors
$U_i\sim N(0,\tau^2I_D)$, before any pretraining. For any even degree-$j$
character put
$\lambda_T(K)=\E_{x,y}\chi_T(x)K(x,y)\chi_T(y)$.
Then
\[
                   \E\lambda_T(K_e^0)\le\tau^2 L_{r,j}. \tag{21}
\]
In particular, if $\mathcal T$ is the family of all degree-$r-2k$
characters and $r-2k\ge4$, with probability at least $1-\xi$ over the
initial features, simultaneously for every $T\in\mathcal T$, every
$q\in\mathcal H(K_e^0)$ with $\|q-\chi_T\|_2\le1/2$ obeys
\[
 \|q\|_{\mathcal H(K_e^0)}^2
       \ge\frac{\xi}{4\tau^2L_{r,r-2k}\binom r{2k}}.    \tag{22}
\]
For fixed $k$ this is exponential in $r$, whereas the pretrained
certificate has polynomial norm and arbitrarily small fixed error.
The probabilities for initial and posterior features refer to their
respective laws, rather than to a specified coupling or training path.

Here is a direct verification of (21). The standard unit-Gaussian ReLU
feature kernel on the cube is
\[
 K_0(x,y)=(r+1)\left\{\frac c4+
 \frac1{2\pi}\left[1+\sum_{l\ge1}a_lc^{2l}\right]\right\},
 \quad c=\frac{1+x\cdot y}{r+1},\quad
 a_l=\frac{\binom{2l}{l}}{4^l(2l-1)^2}\le l^{-5/2}.
\]
Its degree-$j$ coefficient in $c^n$ is the probability that the $j$
specified nonzero symbols have odd multiplicity, and all other nonzero
symbols have even multiplicity, in $n$ independent uniform draws from
$\{0,\ldots,r\}$. This vanishes for $n<j$. Independent Poisson counts
with means $n/(r+1)$ have odd probability at most $1/2$ at each specified
symbol. Conditioning their total to be $n$, whose probability is at
least $1/(e\sqrt n)$, bounds the coefficient by $e\sqrt n\,2^{-j}$.
Summing $a_l\sqrt{2l}$ from $l\ge j/2$ and using
$\sum_{l\ge j/2}l^{-2}\le1/(j/2-1)$ proves the stated bound.
Scaling incoming weights by $\tau$ scales the kernel by $\tau^2$.
Even projection preserves every even Fourier coefficient, proving (21).

A union of Markov bounds over $\binom r{2k}$ characters gives the common
event $\lambda_T(K_e^0)\le\tau^2L_{r,j}\binom r{2k}/\xi$.
For any RKHS function,
\[
 |\E q\chi_T|\le\|q\|_{\mathcal H(K)}\sqrt{\lambda_T(K)}.
\]
An $L^2$ approximation error at most $1/2$ makes the left side at least
$1/2$, proving (22). Conversely, the trained comparator gives
$\lambda_T(K_e)\ge(1-\Delta_A)^2/B_A^2$ whenever $\Delta_A<1$.
This proves an exponential-to-polynomial improvement in the relevant
feature correlations as well as the approximate-predictor norm.
It is not a lower bound for every algorithm that is given the original
input coordinates, or for arbitrary nonlinear use of the frozen features.


\section{An alternative using a linear input skip}
Even feature evaluation can be replaced by an explicit linear skip.
This is a specified augmentation of the downstream architecture, not a
claim about the unaugmented frozen kernel. Suppose $E\le B$ with
$B>0$, and use
\[
 K_{\rm skip}(x,y)=\frac1N\sum_i h_{U_i}(x)h_{U_i}(y)
                                      +\frac B8 x\cdot y.
\]
Under the same hypotheses, there is a comparator with the same error
bound $\Delta_A$ and
\[
 \|q_A^{\rm skip}\|_{\mathcal H(K_{\rm skip})}^2
       \le4B\left(\frac{\rho}{s b_j}\right)^2,
 \qquad \sup_xK_{\rm skip}(x,x)\le2BD+Br/8.
 \tag{23}
\]
To prove this, let $F_c^A$ be the projected, modulated critical network
from the proof. Its odd part is exactly $\ell_A\cdot x$, with
\[
 \ell_A=\frac1{2N}\sum_i a_iR_i\chi_A(z_i)z_i/\sqrt r,
 \qquad \|\ell_A\|\le E/2.
\]
Set $q_A^{\rm skip}=b_r(F^A-\ell_A\cdot x)/(s b_j)$.
The error from replacing $F^A$ by $F_c^A$ is still at most $v_0$;
after subtracting its odd linear part the remainder is $P_eF_c^A$.
The same triangle inequality proves (4). The original-feature
coefficients cost at most $2E(\rho/(s b_j))^2$, and the skip
coefficients cost at most
$2E^2B^{-1}(\rho/(s b_j))^2$. Their sum is bounded by (23).
The radius estimate follows directly.

A downstream convex readout fit does not need to know which coordinates
were deleted: these are comparator-existence bounds for the augmented
feature space. For target degree at least two, a linear skip has zero
Walsh quadratic form on the target. Consequently it does not change the
initial-feature lower bound (22). The learned-versus-initial separation
and the polynomial downstream orders therefore remain valid for this
single-evaluation architecture with an input skip.

\section{An exact obstruction for the unmodified raw features}
The downstream transformation cannot be dropped from a deterministic
near-minimum-energy theorem. The following example attains the sharp
source energy, yet its raw frozen feature space has almost unit error
on a parity only two coordinates away. It also has polynomial-width,
arbitrarily near-minimum-energy versions.

\begin{theorem}[Raw-feature obstruction at optimal source energy]
Let $T\subset[r]$ have even cardinality $2\le j\le r-2$, with $1\in T$.
There is a width-$2^{r-1}$ network with $F=s\chi$, $E=s/\rho$, and
raw frozen feature space $\mathcal V$ satisfying
\[
 \inf_{g\in\mathcal V}\|g-\chi_T\|_2^2
                  =\frac1{1+4r b_j^2}.                 \tag{24}
\]
For $j=r-2k$ with fixed $k\ge1$, the right side tends to one and is
$1-O_k(r^{-2k-1})$.

For every $0<\varepsilon\le1$ and $0<\zeta<1$, there is also a
random construction with exactly $E=(1+\varepsilon)s/\rho$ and any width
\[
 N\ge\frac{4(1+\varepsilon)^2}{\rho^2\varepsilon^2}
                                  \log(2^r/\zeta),      \tag{25}
\]
such that, with probability at least $1-\zeta$, $\chi F\ge s$ on the
entire cube. Every realization of its raw feature space satisfies the
lower bound in (24). Thus a polynomial width and an inverse-polynomial
energy excess do not remove this obstruction.
\end{theorem}
\begin{proof}
Set $B=T\setminus\{1\}$, an odd set of size at most $r-3$, and choose
one orientation from each antipodal pair by
\[
             \mathcal Z=\{z\in\{-1,1\}^r:\chi_B(z)=1\}.
\]
Its uniform law has density $1+\chi_B$ relative to the uniform cube.
Put $h_z(x)=\relu(z\cdot x/\sqrt r)$.
For every such orientation,
\[
 \langle h_z,\chi_T\rangle=b_j\chi_T(z)=b_jz_1,
 \qquad \langle h_z,x_1\rangle=z_1/(2\sqrt r).
\]
Consequently every linear combination $g$ obeys
$\langle g,x_1\rangle=\alpha\langle g,\chi_T\rangle$,
where $\alpha=(2\sqrt r\,b_j)^{-1}$.
Orthogonality of $x_1$ and $\chi_T$ gives
\[
 \|g-\chi_T\|_2^2\ge
 (\langle g,\chi_T\rangle-1)^2+
                   \alpha^2\langle g,\chi_T\rangle^2
 \ge\alpha^2/(1+\alpha^2).
\]
The even coefficients of
$q=b_j^{-1}\E_{z\in\mathcal Z}\chi_T(z)h_z$ are exactly those
of $\chi_T$, and its only odd coefficient is $\alpha x_1$.
Indeed even characters have unchanged expectations under $1+\chi_B$;
the degree-one coefficient at $l$ is nonzero only if
$T\mathbin\triangle\{l\}=B$, equivalently $l=1$.
Thus $q=\chi_T+\alpha x_1$, and $q/(1+\alpha^2)$ attains (24).

Similarly
$\E_{z\in\mathcal Z}\chi(z)h_z=b_r\chi$: all even coefficients
except full parity vanish, and an odd coefficient would require
$[r]\mathbin\triangle\{l\}=B$, impossible because $|B|\le r-3$.
Choose $C=\sqrt{s/\rho}$, $U_z=C(0,z/\sqrt r)$, and
$a_z=C\operatorname{sgn}(b_r)\chi(z)$.
The resulting mean-field network is exactly $s\chi$, has energy $s/\rho$,
and has the same raw feature span as the $h_z$. Formula (3) gives
the asserted fixed-codimension order.

For the finite random version draw the $z_i$ independently from
$\mathcal Z$ and replace $C$ by $\sqrt{(1+\varepsilon)s/\rho}$.
Energy is deterministic, and the mean output is $(1+\varepsilon)s\chi$.
For each input the neuron output has second moment at most
$[(1+\varepsilon)s/\rho]^2$, since
$\E_{\mathcal Z}(z\cdot x)^2=r$: the conditioning character is odd
and all degree-two expectations remain zero.
Its absolute value is at most $(1+\varepsilon)s\sqrt r/\rho$.
Bernstein's inequality and a union over $2^r$ inputs show that (25)
bounds every one-sided deviation by $\varepsilon s$ with failure at
most $\zeta$. Here one uses $\rho\sqrt r\le1/2$ and
$\varepsilon\le1$ to absorb the bounded-summand term in the displayed
constant. Each sampled feature still obeys the exact linear relation,
so the obstruction holds for every realization, independently of the
margin event.
\end{proof}

The same construction also separates raw transfer from regularized
square-risk optimality. With $C^2=1/(2\rho)$ and independent orientations
from $\mathcal Z$, its energy is $1/(2\rho)$, its expected output is
$\chi/2$, and
\[
 \E_z\{\rho E+\E_x(1-\chi F)^2-3/4\}
                         \le\frac1{4\rho^2N}.
\]
Thus, for every $\eta>0$ and $N\ge1/(4\rho^2\eta)$, a realization
has population objective at most $3/4+\eta$ while retaining the same
raw-feature obstruction. This conclusion concerns the objective used in
the square-likelihood companion, without a posterior-probability claim.

This is a statement about a possible learned representation, not its
probability under the Gaussian posterior. In particular it does not
contradict the full-data, width-first raw-feature transfer theorem,
where the ordinary posterior background supplies additional features.
It shows that source interpolation, polynomial width, and almost optimal
energy alone cannot justify unmodified frozen-feature transfer.
Even projection or the linear input skip removes the explicit offending
linear component.

There is also a qualitative classification obstruction. Formula (3)
gives $|2\sqrt r\,b_j|<1$ for $r\ge6$ and the indicated $j$.
Thus the annihilating vector
$v(x)=\chi_T(x)-2\sqrt r\,b_jx_1$ has the same sign as $\chi_T$ at
every input. No function in the raw span, even with an added constant
readout bias, can classify $\chi_T$ correctly on the whole cube: such a
function would have strictly positive inner product with $v$. This
argument also covers a fixed constant tie convention. An adversary can
therefore force a mistake at every step from any deterministic learner
whose classifier is always a raw linear readout, by choosing a currently
misclassified input. This is a frozen-representation statement, not an
end-to-end Bayesian mistake lower bound.

For the full-orientation construction, the uniform population square-loss
fit has an even stronger classification failure. In the orthonormal basis
$(\chi_T,x_1)$ the raw kernel's relevant block, before its common positive
scale factor $C^2$, is
\[
 \begin{pmatrix}b_j^2&b_j/(2\sqrt r)\\
                 b_j/(2\sqrt r)&1/(4r)\end{pmatrix}.
\]
This two-dimensional subspace is invariant under the kernel operator:
the orientation law $1+\chi_B$ pairs each degree-one character with only
one even character. Consequently population ridge regression, at every
finite nonnegative regularization parameter $\gamma$, predicts a positive
multiple of
$b_j^2\chi_T+[b_j/(2\sqrt r)]x_1$.
For $\gamma=0$ this denotes the minimum-square-loss predictor in the raw
feature space. Since $|b_j|<1/(2\sqrt r)$, its classifier is exactly
$\operatorname{sgn}(b_j)x_1$, with uniform error $1/2$.
The same formula holds for a fresh Gaussian readout posterior mean after
any positive equal number of observations of every cube input, with the
corresponding positive ridge parameter. This classification conclusion
uses the explicit invariant block, not a deduction from an $L^2$ norm
lower bound alone.

Equivalently, define the full-support input law
\[
 P_T(x)=2^{-r}\{1-2\sqrt r\,b_j\chi_T(x)x_1\}.
\]
Then $\E_{P_T}\chi_Tg=0$ for every raw linear readout $g$, including
an intercept. Its population square-loss minimizer is the zero function;
its logistic loss is at least $\log2$ by Jensen. The total variation
distance between $P_T$ and uniform is $\sqrt r\,|b_j|$, which is
$\Theta_k(r^{-k-1/2})$ at fixed codimension. This is an explicit small
change of input law under which the frozen representation has no linear
target correlation.

The obstruction is stable under perturbations. Put
$v=\chi_T-2\sqrt r\,b_jx_1$, so that
$\langle h_z,v\rangle=0$ and $\|v\|_2^2=1+4r b_j^2$.
If perturbed features $\widetilde h_i$ satisfy
$N^{-1}\sum_i\|\widetilde h_i-C h_{z_i}\|_2^2\le e^2$, then
Cauchy--Schwarz gives, for every $g$ in their normalized feature RKHS,
\[
 \|g-\chi_T\|_2\le a<\|v\|_2^{-1}
 \quad\Longrightarrow\quad
 \|g\|_{\mathcal H}^2\ge
        \frac{(1-a\|v\|_2)^2}{e^2\|v\|_2^2}.           \tag{26}
\]
At $e=0$ this means that no such approximation exists.
Thus arbitrarily small perturbations restore neither a uniform bounded
readout norm nor a deterministic raw-feature analogue of (4)--(5).


\section{Scope and remaining questions}
The deterministic certificate and the finite-temperature posterior
argument avoid all full-data condensation and conditional local-limit
theorems. They apply at polynomial widths even when exact minimum-energy
representation needs exponential width. The witness is only approximately
energy minimizing, with a controlled inverse-polynomial excess.

The result still concerns exact Bayesian sampling, not an efficient
pretraining algorithm. Its polynomial exponents are conservative. The
prior must be cold enough for the desired transfer accuracy; the argument
does not establish the same result at unit Gaussian scale. It also does
not prove a transfer guarantee for the unsymmetrized frozen kernel alone.

The even feature transformation is essential to this proof. Reweighting
critical ridges can create a linear component that is invisible in the
even Fourier equations. Even projection removes it exactly. The linear-skip construction above supplies one alternative. A further
possibility is to prove that the ordinary posterior background can cancel
it, as happens in the full-cube, width-first posterior calculation.

\section*{Proof of the inherited minimum-coordinate inequality}
For a coordinatewise nondecreasing $g:\{-1,1\}^q\to[0,1]$, permutation
averaging leaves the top Fourier coefficient unchanged. The averaged
function is a nondecreasing sequence on Hamming levels. It is a convex
combination of constants and upper-level thresholds. The alternating
binomial sum of a threshold starting at level $l$ has absolute value
$\binom{q-1}{l-1}$, so its normalized top coefficient is at most
$2^{-q}\binom{q-1}{\lfloor(q-1)/2\rfloor}$.
Flip input signs to make all ridge weights nonnegative. Integrating
coordinate $i$ gives $w_i$ times the top coefficient, on the remaining
$r-1$ coordinates, of
\[
 t\longmapsto\frac{\relu(t+w_i)-\relu(t-w_i)}{2w_i}.
\]
This is a nondecreasing function with values in $[0,1]$.
The preceding bound gives
$|G(u)|\le w_i2^{-(r-1)}\binom{r-2}{(r-2)/2}$.
A zero coordinate has zero full-parity coefficient directly. Minimizing
over $i$ proves (7), including arbitrary bias.

The posterior and transfer theorems are conventional proofs checked in
the same research session, without independent review. The accompanying
Lean files verify the actual single-neuron direction projection, its
finite-network weighted extension, the abstract bounded-modulation
error estimate, the abstract annihilator obstruction with perturbations,
and a concrete six-bit raw-feature example with exact optimal squared
error $64/65$. They do not formalize the Fourier coefficient-ratio
calculation, the posterior concentration argument, or the entire main theorem.
The sharp atom inequality is inherited from
\emph{Minimum Gaussian energy for parity requires exponential ReLU width}
in the existing repository; the finite-width stability and transfer
argument above is the additional result of this note.
\end{document}
